\uuid{v0Qr}
\exo7id{6740}
\titre{exo7 6740}
\auteur{queffelec}
\organisation{exo7}
\datecreate{2011-10-16}
\isIndication{false}
\isCorrection{false}
\chapitre{Autre}
\sousChapitre{Autre}
\cpge{mpsi-mp2i.nombres-complexes.b, pcsi.nombres-complexes.b, ptsi.nombres-complexes.b, mpsi-mp2i.derivabilite.e, pcsi.derivabilite.f, ptsi.derivabilite.e, mpsi-mp2i.integration.c, pcsi.integration.b, ptsi.integration.b, bcpst1.nombres-complexes.b, bcpst1.integration-segment.b, mp-mpi.suites-series-fonctions.a, pc.suites-series-fonctions.a, psi.suites-series-fonctions.a, mp-mpi.series-entieres.b, pc.series-entieres.d, psi.series-entieres.d, pt.series-entieres.d}
\module{Analyse complexe}
\niveau{L3}
\difficulte{}

\contenu{
\texte{
Soit $w=|w|e^{i\theta}$ un nombre complexe.
}
\begin{enumerate}
    \item \question{Montrer que $e^w-1=\int_{[0,w]}e^z\ dz=\int_0^{|w|}
e^{te^{i\theta}}e^{i\theta}\ dt$; montrer ainsi l'inégalité
$|e^w-1|\leq|w|e^{|w|}$ (autre démonstration ?).}
    \item \question{Application. On considère $K$ un compact du plan inclus dans $\Cc^*$ et la suite de fonctions $v_n$
définie sur $K$ par $v_n(z)={1\over z(1+{1\over n})^z }$. Montrer que $v_n(z)$
tend vers ${1\over z}$ uniformément sur $K$.}
\end{enumerate}
}
